\begin{proof}[Beweis des Satzes von Jordan für Kreise\label{jordan:beweis:jordan_kreis}]
    Sei $\gamma(t) = (r \cos(2\pi t), r \sin(2\pi t))$ eine einfache geschlossene Kurve in der Ebene $\mathbb{R}^2$. 
    Das innere Gebiet $I$ ist definiert als die Menge aller Punkte $(x,y)$, die die Ungleichung $x^2 + y^2 < r^2$ erfüllen. 
    Das äussere Gebiet $A$ ist definiert als die Menge aller Punkte $(x,y)$, die die Ungleichung $x^2 + y^2 > r^2$ erfüllen. 
    Die Kurve $\gamma$ selbst bildet den gemeinsamen Rand dieser beiden Gebiete, da sie die Gleichung $x^2 + y^2 = r^2$ erfüllt.    
    Daher hält der Satz von Jordan für Kreise.
\end{proof}